\chapter{Cross-Modal Embedding Architectures}
\label{ch:59}

This chapter establishes the unified embedding framework that integrates
visual, auditory, symbolic, geometric, and narrative modalities into a single
harmonic RSVP--Polyxan manifold.  
The goal is to define a functorial embedding system
\[
\mathsf{Embed} : \mathsf{Modality} \to \mathsf{RSVP\text{-}Man}
\]
whose outputs satisfy curvature-matching, entropy-consistency, and
quine-stability.

\section{Modalities and Their Structures}
\label{sec:59:modalities}

We consider the following modalities:
\[
\mathsf{Modality} 
= \{ \mathrm{Vision},\ \mathrm{Audio},\ \mathrm{Symbolic},\ 
     \mathrm{Geometry},\ \mathrm{Narrative} \}.
\]

\begin{definition}[Modality Structure]
Each modality $M$ is equipped with:
\begin{enumerate}
  \item a data space $D_M$,
  \item a structural operator $T_M : D_M \to D_M$,
  \item a primitive curvature operator $\mathcal{H}_M : D_M \to \mathbb{Z}$,
  \item a tag-sheaf $\mathcal{T}_M$ for semantic refinement.
\end{enumerate}
\end{definition}

\section{The Cross-Modal Embedding Functor}
\label{sec:59:functor}

\begin{definition}[Cross-Modal Functor]
The embedding functor is defined by
\[
\mathsf{Embed}(M)
:= \iota_M : D_M \hookrightarrow \mathcal{M},
\]
where $\mathcal{M}$ is the universal RSVP manifold.
\end{definition}

Functoriality requires:
\[
\iota_{N} \circ F = \iota_{M}
\]
for every modality translation $F : D_M \to D_N$.

\begin{theorem}[Functorial Curvature Matching]
For every modality $M$,
\[
\mathcal{R} \circ \iota_M = \mathcal{H}_M.
\]
\end{theorem}

This ensures that multimodal data embeds into the same geometric curvature
strata.

\section{Entropy-Consistent Multimodal Fusion}
\label{sec:59:entropy}

\begin{definition}[Entropy Profile]
The entropy profile of modality $M$ is
\[
S_M(x) := \mathrm{uncertainty}(T_M(x)).
\]
\end{definition}

\begin{definition}[Cross-Modal Fusion]
Given embeddings $\iota_M$ for each modality,
the fused embedding is the entropy-weighted barycentric combination
\[
\iota_{\mathrm{fuse}}(x_1,\dots,x_k)
= \mathrm{argmin}_{y\in\mathcal{M}}
  \sum_{i=1}^k w_i\, d_{\mathcal{M}}(y,\iota_{M_i}(x_i))^2,
\]
where $w_i = S_{M_i}(x_i)^{-1}$.
\end{definition}

\begin{theorem}[Entropy Consistency]
The fused embedding satisfies
\[
S(y) = \min_i S_{M_i}(x_i).
\]
\end{theorem}

This ensures the combined representation never increases epistemic uncertainty.

\section{Quine-Stable Cross-Modal Alignment}
\label{sec:59:quine}

\begin{definition}[Cross-Modal Quine Frame]
The quine-frame of a multimodal configuration is the smallest Polyxan
hypergraph whose tags match all $\mathcal{T}_{M_i}$.
\end{definition}

\begin{theorem}[Quine-Stable Alignment]
Let $(x_1,\ldots,x_k)$ be multimodal inputs.  
Then the fused embedding $\iota_{\mathrm{fuse}}$ is quine-stable:
\[
\mathcal{G}_{t+1} 
\equiv \mathcal{Q}^m(\mathcal{G}_t)
\quad\text{with } m\le k.
\]
\end{theorem}

Thus multimodality cannot break quine consistency.

\section{Cross-Modal Curvature Synchronisation}
\label{sec:59:synch}

\begin{definition}[Synchronisation Map]
The synchronisation map is
\[
\mathsf{Sync} :
 \prod_{M} \iota_M(D_M) \to \mathcal{M}
\]
defined by curvature projection:
\[
\mathsf{Sync}(x_M) = 
\mathrm{proj}_{\mathcal{R}}
   \bigl(\{\iota_M(x_M)\}_{M}\bigr).
\]
\end{definition}

\begin{theorem}[Curvature Synchronisation]
For any multimodal configuration,
\[
\mathcal{R}(\mathsf{Sync}(x_M))
=
\min_M \mathcal{R}(\iota_M(x_M)).
\]
\end{theorem}

The fused point always descends to the lowest available curvature stratum.

\section{Cross-Modal Embeddings as a Monoidal Category}
\label{sec:59:monoidal}

\begin{definition}[Cross-Modal Monoidal Product]
For modalities $M,N$ define
\[
M \otimes N := \mathrm{Fuse}(M,N),
\]
with unit object the silent modality $\mathsf{Silence}$.
\end{definition}

\begin{theorem}[Monoidal Curvature Law]
\[
\mathcal{H}_{M\otimes N}
= \min(\mathcal{H}_M,\mathcal{H}_N),
\]
and similarly for entropy.
\end{theorem}

Thus multimodal fusion is canonically monoidal.

\section{Universal Cross-Modal Embedding Theorem}
\label{sec:59:universal}

\begin{theorem}[Universal Cross-Modal Embedding]
There exists a unique embedding
\[
\iota_{\mathrm{univ}} : \mathsf{Modality}^\star \to \mathcal{M}
\]
such that:
\begin{enumerate}
  \item it factorises all $\iota_M$,
  \item it preserves curvature depth,
  \item it preserves entropy profiles,
  \item it is quine-stable,
  \item it is initial in the monoidal category of embeddings.
\end{enumerate}
\end{theorem}

\begin{corollary}[Universality of Multimodal Meaning]
All multimodal semantic structures are uniquely embedded into the same 
RSVP manifold, and interact through curvature descent.

\end{corollary}

\section{Conclusion}
\label{sec:59:conclusion}

We have unified vision, audio, symbolic structure, geometry, and narrative into
a single, curvature-consistent, entropy-consistent, quine-stable cross-modal
embedding architecture.

This chapter completes the multimodal integration layer.  
The next chapter, Chapter~\ref{ch:60}, builds the operating system that schedules
all RSVP and Polyxan processes across modalities.

